\chapter{Divisions de marché des banques}\label{ch:in:bank-markets-divisions}

Chaque grande banque publie le revenu de ses activités de marché sur deux lignes~: les produits de taux, les changes et les matières premières (\omterm{def:in:bank-markets-divisions:ficc}{FICC}), et les actions.
La répartition varie d'une banque à l'autre plus que leur taille. En 2025, une banque américaine a tiré 64~\% de son revenu de marché en
actions, et une autre en tire 26~\%~; une banque française a tiré 61~\% en actions, tout en décrivant son activité de gros comme disposant de
un «~leadership mondial distinctif dans les dérivés actions~»~; une banque allemande n'en tire aucun, ayant annoncé en juillet 2019
qu'elle se retirerait de la vente et de la négociation d'actions. La composition du revenu reflète la force de la banque en produits, et elle détermine les postes quant que la banque offre~:
une banque forte sur les dérivés actions emploie des quants de pricing pour les produits structurés et la volatilité, une
banque forte sur les taux emploie des quants de courbe et de XVA. Ce chapitre lit les divisions de marché des banques comme employeurs, à travers ce qu'elles
publient.

\section{Ce que les banques publient sur leurs activités de marché}

Les banques sont les employeurs les plus transparents de ce livre. Les banques cotées déposent des rapports annuels avec les résultats par secteur
(livre~16, chapitre~5), et, au sein du secteur de la banque d'investissement, la plupart publient le revenu de marché par ligne de produits, trimestre par
trimestre. Ce qu'elles publient, et ce qu'elles ne publient pas, conditionne ce qu'un candidat peut apprendre.

\begin{definition}[FICC]\label{def:in:bank-markets-divisions:ficc}
Le \emph{FICC}\index{FICC} (fixed income, currencies and commodities) est l'activité de marché d'une banque qui négocie
des obligations d'État et d'entreprises, des dérivés de taux et de crédit, des changes, des matières premières et leur
financement~; sa ligne symétrique est celle des actions~: actions au comptant, dérivés actions et services de prime brokerage.
\end{definition}

Chaque banque trace la frontière à sa façon. L'une publie «~fixed income markets~» et «~equity markets~»~; une autre «~\omterm{def:in:bank-markets-divisions:ficc}{FICC}~» et
«~equities~» au sein d'un secteur qui comprend aussi le conseil~; une troisième publie «~fixed income~» et «~equity~» au sein de son
secteur «~institutional securities~»~; une banque européenne publie «~Equity \& Prime Services~»~; une autre inclut les
activités de financement dans les produits de taux. Les comparaisons portent donc sur des parts au sein de chaque banque, et non sur des montants d'une banque à l'autre, et
les intitulés propres à chaque banque sont conservés à côté de ses chiffres.

\begin{omfigure}
\begin{tikzpicture}[font=\footnotesize,>=stealth,
  top/.style={draw,rounded corners,fill=omDef!12,minimum width=3.4cm,minimum height=0.6cm,align=center},
  desk/.style={draw,rounded corners,fill=omProp!8,minimum width=4.6cm,minimum height=0.72cm,align=center}]
\node[top] (m) at (4.5,0) {division de marché};
\node[top] (f) at (1.5,-1.2) {FICC};
\node[top] (e) at (7.5,-1.2) {actions};
\draw (m.south) -- ++(0,-0.3) -| (f.north); \draw (m.south) -- ++(0,-0.3) -| (e.north);
\foreach \n/\y/\t/\q in {r/-2.2/taux/{quants courbe et options de taux},c/-3.15/crédit/{quants crédit et XVA},
  x/-4.1/change/{quants options de change et e-FX},k/-5.05/matières premières/{quants modèles de matières premières}}
  {\node[desk] (\n) at (1.5,\y) {\t\\[-1pt]{\scriptsize\itshape \q}};}
\foreach \n/\y/\t/\q in {q/-2.2/actions au comptant/{quants exécution et indices},d/-3.15/dérivés actions/{quants volatilité et structuration},
  p/-4.1/prime services/{quants financement et marge}}
  {\node[desk] (\n) at (7.5,\y) {\t\\[-1pt]{\scriptsize\itshape \q}};}
\draw (f.south) -- (r.north); \draw (e.south) -- (q.north);
\end{tikzpicture}

\omcaption{Les deux lignes de marché que publient les banques et les desks qui les composent, avec les emplois quantitatifs que chaque desk tend à
employer (le chapitre~18 décrit les rôles de quants en banque). Les banques tracent ces lignes différemment~; voici la forme
commune.}\label{fig:in:bank-markets-divisions:tree}
\end{omfigure}

\section{Banques américaines}

Les cinq plus grandes banques américaines publient leurs lignes de marché dans leurs rapports annuels sur le Form 10-K, trois années à la fois. En
2025, leurs parts en actions allaient de 26~\% à 64~\%
(la \cref{fig:in:bank-markets-divisions:mix}), et la part a augmenté dans chacune des quatre banques qui publient des chiffres 2023, de 4.5 à
7.7 points de pourcentage, les revenus actions ayant progressé plus vite que ceux des taux. Les deux banques dont les activités relèvent
de la banque d'investissement et des marchés plutôt que de la banque de détail tirent plus de la moitié de leur revenu de marché en
actions~; la plus grande banque de détail et d'entreprises en tire environ trois dollars sur huit.

\begin{dated}{2025-12}{Revenus de marché de neuf banques par ligne de produits}\label{dat:in:bank-markets-divisions:mix}
\footnotesize
\begin{tabular}{@{}lllrrr@{}}
\toprule
banque & lignes publiées & devise & \omterm{def:in:bank-markets-divisions:ficc}{FICC} & actions & part des actions\\
\midrule
Morgan Stanley & Fixed Income~; Equity & M\$ & 8\,716 & 15\,631 & 64.2~\%\\
Societe Generale & Fixed Income and Currencies~; Equities & M\euro & 2\,346 & 3\,634 & 60.8~\%\\
Goldman Sachs & \omterm{def:in:bank-markets-divisions:ficc}{FICC}~; Equities & M\$ & 14\,522 & 16\,535 & 53.2~\%\\
BNP Paribas & \omterm{def:in:bank-markets-divisions:ficc}{FICC}~; Equity \& Prime Services & M\euro & 5\,522 & 4\,046 & 42.3~\%\\
Bank of America & \omterm{def:in:bank-markets-divisions:ficc}{FICC}~; Equities (sales and trading) & M\$ & 12\,267 & 8\,604 & 41.2~\%\\
Barclays & \omterm{def:in:bank-markets-divisions:ficc}{FICC}~; Equities & M\pounds & 5\,429 & 3\,225 & 37.3~\%\\
JPMorgan Chase & Fixed Income Markets~; Equity Markets & M\$ & 22\,532 & 13\,250 & 37.0~\%\\
Citigroup & Fixed Income Markets~; Equity Markets & M\$ & 16\,226 & 5\,744 & 26.1~\%\\
Deutsche Bank & Fixed Income \& Currencies (incl.\ financing) & M\euro & 9\,600 & -- & 0~\%\\
\bottomrule
\end{tabular}

\smallskip\noindent Chiffres annuels 2025 issus des Forms 10-K et des communiqués de résultats des banques, classés par part des actions.
Les banques américaines publient aussi les chiffres 2023~: les parts des actions s'élevaient à 31.4~\% (JPMorgan Chase), 48.7~\% (Goldman Sachs), 56.5~\%
(Morgan Stanley) et 21.6~\% (Citigroup).
\end{dated}

\section{Les banques françaises et les dérivés actions}

Deux banques françaises se situent de part et d'autre du milieu~: l'une tire 61~\% de son revenu de marché en actions, l'autre 42~\%.
La première décrit sa banque de gros comme offrant «~un leadership mondial distinctif dans les dérivés actions, le financement structuré
et l'ESG~». Au sens où on l'entend ici, les dérivés actions sont pour l'essentiel des produits structurés (livre~5, chapitre~19)~: des notes structurées
dont le remboursement dépend d'indices boursiers et d'actions individuelles, vendues par des réseaux à des épargnants et à des investisseurs institutionnels, et couvertes par les
desks de trading de la banque (livre~9, chapitre~28). Une activité de ce type emploie beaucoup de quants par dollar de revenu~:
structureurs conçoivent les produits, les quants de bibliothèque construisent les pricers, les stratégistes de desk couvrent les livres exotiques, les validateurs
vérifient les modèles (chapitre~18).

\begin{remark}[Pourquoi la spécialisation persiste]\label{rem:in:bank-markets-divisions:why}
La force d'une banque en produits est persistante. Les bibliothèques de pricing, le personnel formé, les relations de distribution et le risque du
portefeuille existant demandent des années à construire et ne se déplacent pas vite~; une banque forte sur un produit tend à le rester, et celle
qui quitte un produit y revient rarement. Les chiffres publiés montrent la composition, non ses causes~; les causes que donnent les banques figurent
dans leurs propres descriptions, qui sont des énoncés de stratégie, et non des preuves.
\end{remark}

\section{Autres banques européennes}

Une banque britannique tire 37~\% de son revenu de marché en actions, et une banque allemande n'en tire aucun~: ayant annoncé en juillet 2019
qu'elle allait «~sortir de son activité Equities Sales \& Trading, tout en conservant une activité ciblée sur les marchés de capitaux actions
»~, elle a publié un revenu de taux et de changes de 9.6~milliards d'euros en 2025, y compris les
activités de financement. Une sortie de ce type supprime d'un coup une filière d'emplois quant chez un même employeur, et les personnes rejoignent
les banques et les firmes qui restent. Convertis aux taux de change moyens de 2025 de la Banque centrale européenne, les revenus de marché des neuf banques
ont totalisé 173.9~milliards de dollars, dont 72.7~milliards en actions~; les trois banques les plus exposées aux actions en ont tiré 49.9~\%
de ce revenu actions.

\begin{omfigure}
\begin{tikzpicture}
  \begin{axis}[width=11cm,height=6.2cm,xbar stacked,bar width=9pt,xmin=0,xmax=40,xlabel={\small Revenu de marché 2025 (milliards de \$, taux moyens BCE)},
  ytick={0,1,2,3,4,5,6,7,8},yticklabels={{Morgan Stanley},{Societe Generale},{Goldman Sachs},{BNP Paribas},{Bank of America},{Barclays},{JPMorgan Chase},{Citigroup},{Deutsche Bank}},y dir=reverse,
  tick label style={font=\footnotesize},grid=major,grid style={gray!20},table/col sep=comma,
  legend style={legend cell align=left,font=\scriptsize,at={(0.5,-0.2)},anchor=north,/tikz/every even column/.append style={column sep=8pt}},legend columns=2]
\addplot[fill=omThm!55,draw=omThm] table[y=k,x=equities_bn]{figdata/industry/07-bank-markets-divisions/mix2025.csv};
\addplot[fill=omDef!55,draw=omDef] table[y=k,x=ficc_bn]{figdata/industry/07-bank-markets-divisions/mix2025.csv};
\legend{actions,taux et change (FICC)}
\end{axis}
\end{tikzpicture}

\omcaption{Revenus de marché 2025 des neuf banques, actions et taux et change, convertis en dollars américains aux taux moyens annuels
de la BCE, classés par part des actions (et non par taille). Les lignes de chaque banque lui sont propres et ne sont pas strictement
comparables en montant. Données~: Forms 10-K et communiqués de résultats pour 2025, via \texttt{in\_banks.usd}.}\label{fig:in:bank-markets-divisions:mix}
\end{omfigure}

\section{Classements~: qui les établit et ce qu'ils mesurent}

\begin{definition}[Réservoir de revenus, classement]\label{def:in:bank-markets-divisions:league}
Un \emph{réservoir de revenus}\index{réservoir de revenus} (revenue pool) est le revenu total que les dealers gagnent ensemble sur un produit ou un groupe
de clients sur une période, tel qu'estimé par un analyste. Un \emph{classement}\index{classement} (league table) classe les dealers selon leur
part d'un réservoir de revenus, de volumes ou de l'activité des clients, telle qu'estimée et publiée par une société d'analyse, une
publication sectorielle ou un fournisseur de données, pour une période donnée.
\end{definition}

Les banques citent des \omterm{def:in:bank-markets-divisions:league}{classements} dans leurs présentations, et les candidats les lisent comme un palmarès des employeurs. Ce ne sont ni des
documents déposés ni des comptes audités~: ce sont des estimations établies par des sociétés d'analyse à partir des déclarations propres aux banques et
d'enquêtes auprès des clients. L'une de ces sociétés décrit son service comme une aide permettant aux banques de «~mesurer la croissance du revenu, la part du portefeuille et la fidélisation
à travers les ventes, les produits et les services~», avec pour
clients «~les grandes banques d'investissement corporate et les dealers mondiaux et régionaux~». Le rang d'un \omterm{def:in:bank-markets-divisions:league}{classement} dépend de ses définitions --- quels produits, quelles régions, quelle période --- et de
l'organisme qui transmet ses données. Ce livre ne cite aucun \omterm{def:in:bank-markets-divisions:league}{classement}. Lorsqu'un tableau est utile, il indique son éditeur, sa date
et ce qu'il mesure, et un candidat doit en faire autant.

\section{Effectifs de quants~: ce qui est public et ce qui ne l'est pas}

Aucune banque de ce chapitre ne publie le nombre de quants qu'elle emploie. Les rapports annuels donnent l'effectif total (chapitre~1~: 318\,512 dans une
banque, 47\,400 dans une autre) et parfois le nombre dans la banque de financement et d'investissement, mais pas le
nombre dans la valorisation, le risque ou la validation. Les éléments publics sont indirects~: les déclarations sur les conditions de travail du chapitre~14 montrent
les intitulés de postes quant des banques et leurs fourchettes de rémunération par employeur, et les informations sur la rémunération du chapitre~14 comptent
les preneurs de risques significatifs par domaine d'activité. Un candidat qui veut connaître la taille d'une équipe demande à l'équipe.

\section{Tutoriel~: où se trouvent les dérivés actions}\label{tut:in:bank-markets-divisions:tut}

\begin{tutorial}
\textbf{Objectif.} Tabuler le revenu de marché de neuf banques par ligne de produit à partir de leurs déclarations, le convertir dans une seule devise et
classer les banques par structure. \textbf{État final~:} l'encadré daté et la \cref{fig:in:bank-markets-divisions:mix}.
\begin{tutsteps}
\item \textbf{Les lignes.} \texttt{data/industry/bank\_markets\_revenue.csv} contient une ligne par banque et par année~: les deux
lignes telles que la banque les publie, la devise, la ligne du grand livre qui reprend la déclaration, et les intitulés propres à la banque.
\item \textbf{Les taux.} \texttt{data/industry/ecb\_fx\_annual.csv} contient les taux de référence annuels moyens de la BCE~;
\texttt{firm.bankmix.load\_fx} les convertit en dollars américains par euro et par livre.
\item \textbf{La structure.} \texttt{firm.bankmix.mix\_table(rows, 2025)} classe les banques par part des actions
(\cref{lst:in:bank-markets-divisions:mix})~; \texttt{in\_banks.changes()} donne la variation de la part de chaque banque entre sa
première et sa dernière année publiées.
\omcode{code/firm/bankmix/firm_bankmix.py}{52}{70}{La part des actions, le classement par structure et la part du revenu actions
gagnée par les banques les plus tournées vers les actions.}\label{lst:in:bank-markets-divisions:mix}
\item \textbf{Concentration.} \texttt{in\_banks.top3()} convertit le revenu actions en dollars et calcule la part des
trois banques les plus tournées vers les actions.
\end{tutsteps}
\end{tutorial}

\textbf{Pour aller plus loin.} Ajouter une banque à partir de son communiqué de résultats, avec ses propres intitulés de lignes (exercice~7)~; recalculer
les totaux en dollars aux taux de 2024 pour voir quelle part de la variation annuelle tient à la devise.

\section{Construction~: la structure des activités de marché}\label{bld:in:bank-markets-divisions:bankmix}

\begin{build}
\bfield{Rôle} Mettre le revenu de marché des banques en deux lignes standard et une seule devise, afin que les chapitres~12 et~18 puissent lire la force produit d'une
banque à côté de ses données de rémunération et de ses postes.
\bfield{Interface} \texttt{firm.bankmix}~: \texttt{Row(bank, year, ficc, equities, currency, ledger, lines)}~;
\texttt{load(path)}~; \texttt{load\_fx(path)}~; \texttt{to\_usd(row, fx)}~; \texttt{equities\_share(row)}~;
\texttt{mix\_table(rows, year)}~; \texttt{mix\_change(rows, bank, y0, y1)}~; \texttt{top\_k\_equities\_share(rows, fx, year, k)}.
\bfield{Règles} Chaque ligne conserve les intitulés propres à la banque et sa source~; les montants sont convertis au taux moyen de l'année,
jamais au taux spot~; les banques sont classées par structure, jamais par taille seule ni par un quelconque \omterm{def:in:bank-markets-divisions:league}{classement}.
\bfield{Tests d'acceptation} \texttt{code/firm/bankmix/tests/}~: parts et ordre sur des lignes construites~; variation de la part en points de pourcentage~;
conversion par le cours croisé euro~; part des $k$ premières~; aller-retour CSV.
\bfield{Pour aller plus loin} Lignes trimestrielles et saisonnalité de la structure~; régression de la part actions de chaque banque sur la moyenne annuelle
de l'indice VIX (tableau du livre~16).
\end{build}

\begin{omsources}
\item Forms 10-K pour 2025 de JPMorgan Chase, Goldman Sachs, Morgan Stanley, Citigroup et Bank of America.
\item Communiqués de résultats pour 2025~: BNP Paribas (annexes, 5 février 2026), Societe Generale (6 février 2026), Barclays
(2025 Results Announcement), Deutsche Bank (29 janvier 2026)~; communiqué de presse de Deutsche Bank du 7 juillet 2019.
\item Banque centrale européenne, euro foreign exchange reference rates, annual averages.
\item Crisil Coalition Greenwich, description de ses services de benchmarking.
\end{omsources}

\section{Exercices}

\begin{exercise}[$\star$]\label{exo:in:bank-markets-divisions:1}
Calculer la part des actions dans le revenu de marché de Citigroup en 2023 et en 2025, ainsi que sa variation en points de pourcentage.
\end{exercise}

\begin{exercise}[$\star$]\label{exo:in:bank-markets-divisions:2}
Convertir le revenu \omterm{def:in:bank-markets-divisions:ficc}{FICC} 2025 de la banque britannique, soit 5\,429~millions de livres, en dollars aux taux moyens de la BCE.
\end{exercise}

\begin{exercise}[$\star$]\label{exo:in:bank-markets-divisions:3}
Indiquer quels desks de la \cref{fig:in:bank-markets-divisions:tree} emploieraient le plus de model quants dans une banque
dont la part des actions est de 60~\%, et pourquoi~?
\end{exercise}

\begin{exercise}[$\star\star$]\label{exo:in:bank-markets-divisions:4}
La part des actions de Morgan Stanley est passée de 56.5~\% à 64.2~\% entre 2023 et 2025. Décomposer la variation~: de combien
le revenu actions et le revenu de taux et de change ont-ils chacun augmenté~?
\end{exercise}

\begin{exercise}[$\star\star$]\label{exo:in:bank-markets-divisions:5}
Expliquer pourquoi les montants de la \cref{fig:in:bank-markets-divisions:mix} ne doivent pas être lus comme un \omterm{def:in:bank-markets-divisions:league}{classement} des activités
de marché des banques~?
\end{exercise}

\begin{exercise}[$\star\star$]\label{exo:in:bank-markets-divisions:6}
Un \omterm{def:in:bank-markets-divisions:league}{classement} (league table) place une banque au premier rang dans un produit où elle détient une part de 15~\% du \omterm{def:in:bank-markets-divisions:league}{réservoir de revenus}. Lister quatre questions à poser avant
d'utiliser ce rang pour comparer des employeurs.
\end{exercise}

\begin{exercise}[$\star\star\star$]\label{exo:in:bank-markets-divisions:7}
\emph{Programmation.} Ajouter la banque allemande comme ligne avec zéro action et recalculer la part des revenus actions des trois banques les plus exposées aux actions si la banque allemande
avait conservé une activité actions rapportant 3~milliards d'euros. Qu'est-ce qui change, et pourquoi~?
\end{exercise}

\begin{exercise}[$\star\star\star$]\label{exo:in:bank-markets-divisions:8}
\emph{Trouver la faille.} «~La banque française réalise 61~\% de son revenu en actions et la banque américaine
26~\%~: les quants actions sont mieux payés dans la banque française.~»
\end{exercise}

\section{Problème~: où se trouvent les dérivés actions}

\begin{problem}[{Problème du week-end --- où se trouvent les dérivés actions}]\label{pb:in:bank-markets-divisions:1}
Un quant en dérivés, qui a reçu des offres d'une banque américaine et d'une banque européenne, veut savoir laquelle a l'activité actions la plus solide
et quel poids accorder aux \omterm{def:in:bank-markets-divisions:league}{classements}.

\medskip\noindent\textbf{Partie I --- Les lignes.}
\begin{enumerate}
\item Définir le \omterm{def:in:bank-markets-divisions:ficc}{FICC} et la ligne actions, et nommer les desks de chacune.
\item Pourquoi les comparaisons de parts au sein des banques sont-elles plus sûres que les comparaisons de montants entre elles~?
\item Donner les parts des neuf banques en actions pour 2025, dans l'ordre.
\item Donner les parts des quatre banques américaines en actions pour 2023 et leurs variations jusqu'en 2025.
\item Qu'a annoncé la banque allemande en juillet 2019, et que publie-t-elle aujourd'hui~?
\end{enumerate}

\medskip\noindent\textbf{Partie II --- La conversion.}
\begin{enumerate}[resume]
\item Donner les taux moyens 2025 de la BCE et le nombre implicite de dollars par livre.
\item Convertir en dollars le revenu actions 2025 des banques françaises.
\item Donner en dollars le revenu total de marché et le revenu actions des neuf banques.
\item Donner la part du revenu actions réalisée par les trois banques les plus exposées aux actions.
\item Pourquoi utiliser un taux moyen plutôt que le taux de fin d'année~?
\end{enumerate}

\medskip\noindent\textbf{Partie III --- Les emplois.}
\begin{enumerate}[resume]
\item Quels rôles de quant emploie une banque forte en dérivés actions, et pourquoi en emploie-t-elle beaucoup par dollar de revenu~?
\item Pourquoi la force d'une banque dans un produit est-elle persistante~?
\item Que publient les banques sur leurs effectifs de quants~?
\item Où un candidat peut-il trouver des indices indirects sur les rôles de quant et les rémunérations d'une banque~?
\item Que devient l'emploi de quant dans un produit lorsqu'une banque s'en retire~?
\end{enumerate}

\medskip\noindent\textbf{Partie IV --- Le verdict.}
\begin{enumerate}[resume]
\item Énoncer le \emph{résultat nommé}~: l'étendue des parts en actions des neuf banques en 2025, et la part du
revenu actions réalisée par les trois banques les plus exposées aux actions.
\item Définir un \omterm{def:in:bank-markets-divisions:league}{réservoir de revenus} (revenue pool) et un \omterm{def:in:bank-markets-divisions:league}{classement} (league table).
\item Que doit indiquer un \omterm{def:in:bank-markets-divisions:league}{classement} cité~?
\item Laquelle des deux offres du quant se trouve dans la banque à la composition actions la plus solide, si la banque américaine est celle à 26~\%
et la banque européenne celle à 61~\%~?
\item En deux phrases, que doit peser le quant au-delà de la composition~?
\end{enumerate}
\end{problem}

\section{Questions d'entretien}

\begin{interviewq}[$\star$ \iqroles{banque}]\label{iq:in:bank-markets-divisions:1}
Que signifie \omterm{def:in:bank-markets-divisions:ficc}{FICC}, et en quoi son revenu diffère-t-il en nature du revenu actions~?
\end{interviewq}

\begin{interviewq}[$\star$ \iqroles{banque, trader}]\label{iq:in:bank-markets-divisions:2}
Pourquoi le revenu actions d'une banque pourrait-il croître plus vite que son revenu de taux une année où les cours montent
et où la volatilité augmente~?
\end{interviewq}

\begin{interviewq}[$\star\star$ \iqroles{banque}]\label{iq:in:bank-markets-divisions:3}
Une banque vend des notes structurées à des épargnants particuliers par l'intermédiaire de ses agences. Quels desks et quels quants cette activité demande-t-elle,
de la conception du produit à la couverture~?
\end{interviewq}

\begin{interviewq}[$\star\star$ \iqroles{banque, risque}]\label{iq:in:bank-markets-divisions:4}
Pourquoi une banque abandonnerait-elle totalement la vente et le trading d'actions plutôt que de la réduire~?
\end{interviewq}

\begin{interviewq}[$\star\star$ \iqroles{chercheur}]\label{iq:in:bank-markets-divisions:5}
Vous disposez de quatre années de deux lignes de revenus pour neuf banques. Comment testeriez-vous si la composition est plus stable
que le niveau~?
\end{interviewq}

\begin{interviewq}[$\star\star\star$ \iqroles{banque}]\label{iq:in:bank-markets-divisions:6}
Le revenu de votre desk de dérivés actions a augmenté de 30~\% en un an. Quelle part en attribueriez-vous au marché, quelle part
à la franchise, et comment le distingueriez-vous~?
\end{interviewq}
